\section{Signed receiving bounds at the original quartet degrees}

Retain the original coefficient maps and every constant of ODD1--14.
Thus $m\geq1$, $k\geq3$, $q=(k+1)^2[1+k(m-1)]$,
$p=8m-9/2$, $M=21+4m$, $\eta=(p-1)/p$, and
$z_k=kT_k/q$. The constants $C_{\rm tail},C_s,C_E,C_d,g_0$
and $T_k$ retain exactly their definitions in ODD3--5.
For each of the four degrees $N=q-1,q,2q-1,2q$ put
\[
\begin{aligned}
 t_k^{\rm ODD}&=C_{\rm tail}\min\{2^{1-p},z_k^\eta\}
                 +\frac{(p+1)\log k+C_s}{k},\\
 e_k^{\rm ODD}&=C_Ez_k+C_d\max\{z_k,z_k^\eta\}
                       +\frac{(g_0+2)C_d}{q},\\
 \Omega_{k,N}&=\frac{N+1}{q}t_k^{\rm ODD}
                 +\frac{2MNT_k}{kq}+e_k^{\rm ODD}.
\end{aligned}\tag{ORP1}
\]
These are explicit nonnegative scalars on the unchanged original
source. The finite bound is
\[
 0\leq\Xi_{k,N}\leq kq\Omega_{k,N}\leq kq\Psi_k.
 \tag{ORP2}
\]
Indeed set $R=\max(1,z_k^{-1/p})$ in ODD7. Then
$z_kR=\max\{z_k,z_k^\eta\}$ and
$(1+R)^{1-p}\leq\min\{2^{1-p},z_k^\eta\}$ by ODD8.
Insertion yields the first upper bound of ORP2 without replacing
either $(N+1)/q$ or $N/q$. For the second use
$N/q\leq(N+1)/q\leq2+1/q$ in ORP1, and compare with
the complete expression ODD5 for $\Psi_k$.

Define the two evaluated signed allowances
\[
\begin{aligned}
 D_k^-&=kq(\Omega_{k,q-1}+\Omega_{k,q})\\
 &=kq\left\{\left(2+\frac1q\right)t_k^{\rm ODD}
       +\frac{2M(2q-1)T_k}{kq}+2e_k^{\rm ODD}\right\},\\
 D_k^+&=kq(\Omega_{k,2q-1}+\Omega_{k,2q})\\
 &=kq\left\{\left(4+\frac1q\right)t_k^{\rm ODD}
       +\frac{2M(4q-1)T_k}{kq}+2e_k^{\rm ODD}\right\}.
\end{aligned}\tag{ORP3}
\]
For each of the original quotient, kernel and boundary corrections
$X\in\{\delta_k^\sigma,\Delta F_k^K,\Delta F_k^B\}$ one has
\[
 -D_k^-\leq X\leq D_k^+.
 \tag{ORP4}
\]
To prove this on the actual maps, let $b_N$ be the nonnegative
quotient deficit from HMR41, and $b_N^K,b_N^B$ its two
nonnegative summands from HMR42. The full positive Schur
determinant calculation gives $0\leq b_N^K,b_N^B\leq b_N
\leq\Xi_{k,N}$. For any of these three families the four
original quotient signs yield
$X=-b_{q-1}-b_q+b_{2q-1}+b_{2q}$, with the superscript
inserted for either summand. Dropping the positive two terms
proves $X\geq-\Xi_{k,q-1}-\Xi_{k,q}\geq-D_k^-$.
Dropping the negative two proves
$X\leq\Xi_{k,2q-1}+\Xi_{k,2q}\leq D_k^+$.
The relation kernel, full boundary frame and its determinant
factor are exactly those in HMR39--43; no observation is changed.

For $m\geq2$ set $C_m=p+1+6M=32m+245/2$. If $N/q\to s$
along these original endpoints, then
\[
 \lim_{k\to\infty}\frac{k\Omega_{k,N}}{\log k}=sC_m,
 \qquad
 \limsup_{k\to\infty}\frac{\Xi_{k,N}}{q\log k}\leq sC_m.
 \tag{ORP5}
\]
Here $s=1$ at the first two endpoints and $s=2$ at the last
two. To prove the limit, use $q=(m-1)k^3+O_m(k^2)$,
$T_k/\log k\to3$, and $\eta>1/2$. In particular
\[
 \frac{kz_k^\eta}{\log k}
   =O_h\bigl(k^{1-2\eta}(\log k)^{\eta-1}\bigr)\longrightarrow0,
 \quad
 \frac{kz_k}{\log k}=O_h(k^{-1}),
 \quad \frac{k}{q\log k}\longrightarrow0.
\]
Thus $kt_k^{\rm ODD}/\log k\to p+1$ and
$ke_k^{\rm ODD}/\log k\to0$. The middle term of ORP1,
multiplied by $k/\log k$, tends to $6Ms$. This proves ORP5
and also the exact limits of the explicit allowances:
\[
 \frac{D_k^-}{q\log k}\longrightarrow64m+245,
 \qquad
 \frac{D_k^+}{q\log k}\longrightarrow128m+490.
 \tag{ORP6}
\]
Consequently, for each correction $X$ in ORP4,
\[
 -(64m+245)\leq\liminf\frac{X}{q\log k}
       \leq\limsup\frac{X}{q\log k}\leq128m+490
 \qquad(m\geq2).
 \tag{ORP7}
\]
For $m=1$ the finite ORP1--4 remain valid, and ODD9 gives
$X/(kq)=O_h((\log k/k)^{5/7})$. The limit calculation ORP5
uses the actual cubic growth of $q$ for $m\geq2$; it does
not give a $q\log k$ bound in the simple-zero regime.

On the exact original Gamma return the finite formulas are
\[
\begin{aligned}
 -\mathfrak T_k-D_k^-&\leq\Delta_k^\Gamma
                          \leq-\mathfrak T_k+D_k^+,\\
 \max\{0,F_{\sigma,k}-4q\log(D_hk)-D_k^-\}
       &\leq\mathcal R_k
          \leq F_{\sigma,k}-4q\log(D_hk)+D_k^+.
\end{aligned}\tag{ORP8}
\]
This follows by direct substitution in ODD13, including the
independent proved inequality $\mathcal R_k\geq0$.
In the original two-reference notation replace
$-\mathfrak T_k$ by the complete
$\mathcal R_k^0-\mathcal R_k^\sigma$; they agree wherever
the original signed reference determinant is represented by
$\mathfrak T_k$. Thus ORP7 applies to
$\Delta_k^\Gamma+\mathfrak T_k$ and to
$\mathcal R_k-F_{\sigma,k}+4q\log(D_hk)$ themselves.
It inserts no sign or limiting value for either Gamma residual.

The existing exact TWA15 interval can be retained alongside its
fully evaluated consequence
\[
 \max\{-qE_k^+,-D_k^-\}\leq\delta_k^\sigma
       \leq\min\{qE_k^-,D_k^+\}.
 \tag{ORP9}
\]
The maximum and minimum use only already proved bounds on the
same scalar. ORP9 evaluates the determinant side of the interval;
it does not replace any exactly known smaller deficit or the
separate adjacent monic deficits in TWA6.
